\indent ಷಟ್ + ಮಾಸ = ಷಣ್ + ಮಾಸ = ಷಣ್ಮಾಸ (ಟಕಾರಕ್ಕೆ ಮಕಾರ ಪರವಾಗಿದ್ದು ಣಕಾರ ಆದೇಶವಾಗಿದೆ\par

\indent ಚಿತ್ + ಮಯ = ಚಿನ್ + ಮಯ = ಚಿನ್ಮಯ (ತಕಾರಕ್ಕೆ ಮಕಾರ ಪರವಾಗಿದ್ದು ನಕಾರ ಆದೇಶವಾಗಿದೆ ಇವು ಅನುನಾಸಿಕ ಸಂಧಿಗೆ ಉದಾಹರಣೆಗಳು ; ಅನುನಾಸಿಕ ಸಂಧಿ ಯಾಗುವ ಸಂದರ್ಭಗಳನ್ನು ಕೆಳಗೆ ಕೊಟ್ಟಿರುವ ಸೂತ್ರದಿಂದ ತಿಳಿಯಬಹುದು\par

\begin{kbvrule}
\noindent{\kbvsecfont ಸೂತ್ರ :}\hspace{0.5em}ಪೂರ್ವ ಶಬ್ದದ ಅಂತ್ಯದಲ್ಲಿರುವ ಕ, ಚ, ಟ, ತ, ಪ ವ್ಯಂಜನಗಳಿಗೆ ಯಾವುದೇ ಅನುನಾಸಿಕಾಕ್ಷರ ಪರವಾದರೂ ಕ ಚ ಟ ತ ಪ ಗಳಿಗೆ ಙ, ಞ, ಣ, ನ, ಮ ಗಳು ಕ್ರಮವಾಗಿ ಆದೇಶಗಳಾಗಿ ಬರುವುದಾದರೆ ಅದು ಅನುನಾಸಿಕ ಸಂಧಿಯೆನಿಸುವುದು.
\end{kbvrule}

\kbvexercise{ಅಭ್ಯಾಸ ಪತ್ರಿಕೆ ೭}

\vspace{2.5pt plus 3pt minus 0.5pt}\noindent\hspace*{1.2em}\kbvnum{{[}1{]}}\hspace{0.55em}ಸ್ವರ ಸಂಧಿ ಎಂದರೇನು ? ಸಂಸ್ಕೃತ ಸ್ವರ ಸಂಧಿಗಳಲ್ಲಿ ಎಷ್ಟು ಪ್ರಕಾರಗಳಿವೆ ?\par

\vspace{2.5pt plus 3pt minus 0.5pt}\noindent\hspace*{1.2em}\kbvnum{{[}2{]}}\hspace{0.55em}ವ್ಯಂಜನ ಸಂಧಿ ಎಂದರೇನು ? ಮುಖ್ಯವಾದ ಮೂರು ಸಂಸ್ಕೃತ ವ್ಯಂಜನ ಸಂಧಿಗಳನ್ನು ತಿಳಿಸಿ.\par

\vspace{2.5pt plus 3pt minus 0.5pt}\noindent\hspace*{1.2em}\kbvnum{{[}3{]}}\hspace{0.55em}ಸವರ್ಣ ದೀರ್ಘ ಸಂಧಿ ಎಂದರೇನು? ಒಂದು ಉದಾಹರಣೆ ಸಹಿತ ವಿವರಿಸಿ.\par

\vspace{2.5pt plus 3pt minus 0.5pt}\noindent\hspace*{1.2em}\kbvnum{{[}4{]}}\hspace{0.55em}ಗುಣ ಸಂಧಿ ಎಂದರೇನು ? ಉದಾಹರಣೆಗಳ ಸಹಿತ ತಿಳಿಸಿ.\par

\vspace{2.5pt plus 3pt minus 0.5pt}\noindent\hspace*{1.2em}\kbvnum{{[}5{]}}\hspace{0.55em}ವೃದ್ಧಿ ಸಂಧಿ ಎಂದರೇನು ? ಉದಾಹರಣೆಗಳ ಸಹಿತ ತಿಳಿಸಿ.\par

\vspace{2.5pt plus 3pt minus 0.5pt}\noindent\hspace*{1.2em}\kbvnum{{[}6{]}}\hspace{0.55em}ಯಣ್ ಸಂಧಿಯ ಸೂತ್ರವನ್ನು ನಿರೂಪಿಸಿ, ಉದಾಹರಣೆಗಳ ಮೂಲಕ ಸಮರ್ಥಿಸಿ?\par

\vspace{2.5pt plus 3pt minus 0.5pt}\noindent\hspace*{1.2em}\kbvnum{{[}7{]}}\hspace{0.55em}ಜಶ್ತ್ವ ಸಂಧಿಗೆ ಎರಡು ಉದಾಹರಣೆಗಳನ್ನು ಕೊಡಿ.\par

\vspace{2.5pt plus 3pt minus 0.5pt}\noindent\hspace*{1.2em}\kbvnum{{[}8{]}}\hspace{0.55em}ಶ್ಚುತ್ವ ಸಂಧಿಯ ಸೂತ್ರವನ್ನು ನಿರೂಪಿಸಿ.\par

\vspace{2.5pt plus 3pt minus 0.5pt}\noindent\hspace*{1.2em}\kbvnum{{[}9{]}}\hspace{0.55em}ಅನುನಾಸಿಕ ಸಂಧಿ ಎಂದರೇನು ? ಉದಾಹರಣೆ ಸಹಿತ ತಿಳಿಸಿ.\par

\vspace{2.5pt plus 3pt minus 0.5pt}\noindent\hspace*{1.2em}\kbvnum{{[}10{]}}\hspace{0.55em}ಇವುಗಳನ್ನು ಕೂಡಿಸಿ ಬರೆದು ಸಂಧಿಯನ್ನು ತಿಳಿಸಿರಿ.\par

\indent ಅಷ್ಟ್ + ಮಾಸಂ; ಮಾಸ್ + ಮಾನ; ಮಾಸಂ + ಮಾಸಂ\par
